\section{The Fiat-Shamir Transformation}\label{sec:fiat_shamir}

The Fiat-Shamir transformation turns a public-coin interactive oracle reduction into a
non-interactive reduction by deriving verifier challenges from the statement and the prover
messages already seen. ArkLib contains two related formalizations:

\begin{itemize}
    \item the canonical random-oracle Fiat-Shamir transformation, used as the basic reference
    model; and
    \item the duplex-sponge Fiat-Shamir transformation of Chiesa--Orru~\cite{CO25}, which
    derives challenges by incrementally absorbing encoded prover messages into a sponge and
    squeezing encoded verifier challenges.
\end{itemize}

The canonical construction is the simplest machine-checked interface for security transfers.
The duplex-sponge construction is the grant target: it models the deployed sponge-style API,
including the codec layer and the salted paper variant, while reusing the canonical construction
as the reference experiment in the key-lemma security framework.

\subsection{Canonical Random-Oracle Fiat-Shamir}\label{sec:fiat_shamir_canonical}

\begin{definition}[Canonical Fiat-Shamir Challenge Oracle]
    \label{def:fiat_shamir_challenge_oracle_interface}
    For a protocol specification $\pSpec$ and input statement type $\StmtIn$, the canonical
    Fiat-Shamir oracle has one oracle family member for every verifier challenge round $i$.
    Its query is the statement together with the prover-message prefix before round $i$, and
    its response is a challenge of type $\pSpec.\Challenge(i)$:
    \[
        \mathsf{fsChallengeOracle}\;\StmtIn\;\pSpec
        : \OracleSpec\;\pSpec.\ChallengeIdx .
    \]
    This is the transcript-prefix random-oracle model used by the basic security-transfer
    theorems.
    \lean{ProtocolSpec.fsChallengeOracle}
\end{definition}

\begin{definition}[Canonical Prover and Verifier Transform]
    \label{def:canonical_fiat_shamir_transform}
    Given an interactive reduction $R$, the canonical transform applies the Fiat-Shamir prover
    and verifier transformations componentwise:
    \[
        R.\mathsf{fiatShamir}
        : \NonInteractiveReduction
          \;(\forall i,\pSpec.\Message(i))\;
          (\oSpec + \mathsf{fsChallengeOracle}\;\StmtIn\;\pSpec)\;\cdots .
    \]
    The transformed prover computes all original prover messages in one non-interactive proof.
    The transformed verifier reconstructs the full interactive transcript by querying the
    canonical challenge oracle at each verifier round.
    \lean{Prover.fiatShamir}
    \lean{Verifier.fiatShamir}
    \lean{Reduction.fiatShamir}
    \uses{def:fiat_shamir_challenge_oracle_interface}
\end{definition}

\begin{theorem}[Canonical Completeness and Security Transfers]
    \label{thm:canonical_fiat_shamir_security}
    The canonical transformation has machine-checked closeout theorems for the main security
    transfers used by ArkLib: completeness, state-restoration soundness, state-restoration
    knowledge soundness, and honest-verifier zero-knowledge transfer for the canonical
    implementations.
    \lean{Reduction.fiatShamir_completeness_unroll_discharged}
    \lean{Reduction.fiatShamir_soundness_of_stateRestoration}
    \lean{Reduction.fiatShamir_knowledgeSoundness_of_stateRestoration}
    \lean{Reduction.fiatShamir_isHVZK_canonical}
    \lean{Reduction.fiatShamir_isStatHVZK_canonical}
\end{theorem}

\subsection{Duplex-Sponge Oracle and Codec Layer}\label{sec:fiat_shamir_duplex_oracle}

The duplex-sponge construction replaces the transcript-prefix challenge oracle with a fixed-size
ideal permutation interface. A sponge state has rate $R$, capacity $C$, total state size $N=R+C$,
and unit alphabet $\Sigma$. Prover messages are encoded into vectors over $\Sigma$, absorbed into
the sponge, and verifier challenges are squeezed as vectors over $\Sigma$ and decoded into the
protocol's challenge spaces.

\begin{definition}[Duplex-Sponge Challenge Oracle]
    \label{def:duplex_sponge_challenge_oracle}
    The duplex-sponge challenge oracle is the three-part oracle
    \[
        \mathcal{D}_{\mathfrak{S}}
        =
        (\mathsf{StartType} \to \Sigma^C)
        + \mathsf{PermOracle}(\Sigma^N),
    \]
    whose query flavors are:
    \begin{itemize}
        \item $h(x)$, initializing a sponge capacity vector from a statement-like start value;
        \item $p(s)$, a forward ideal-permutation query on sponge states; and
        \item $p^{-1}(s)$, the inverse-permutation query on the same state space.
    \end{itemize}
    ArkLib exposes these flavors through the reducible query constructors
    $\mathsf{dsHashQuery}$, $\mathsf{dsPermQuery}$, and $\mathsf{dsPermInvQuery}$.
    \lean{OracleSpec.duplexSpongeChallengeOracle}
    \lean{OracleSpec.dsHashQuery}
    \lean{OracleSpec.dsPermQuery}
    \lean{OracleSpec.dsPermInvQuery}
\end{definition}

\begin{definition}[Protocol Codec]
    \label{def:duplex_sponge_codec}
    A codec for a protocol $\pSpec$ over the sponge alphabet $\Sigma$ supplies:
    \begin{itemize}
        \item message sizes and challenge sizes;
        \item encoders $\phi_i : \pSpec.\Message(i) \to \Sigma^{\ell_P(i)}$ with injectivity;
        \item decoders $\psi_i : \Sigma^{\ell_V(i)} \to \pSpec.\Challenge(i)$;
        \item a decoding-bias bound $\epsilon_{\mathsf{cdc},i}$ for each challenge round; and
        \item a sampler for challenge preimages under $\psi_i$.
    \end{itemize}
    The codec is the formal bridge between ArkLib protocol messages and the fixed alphabet used
    by the duplex sponge.
    \lean{ProtocolSpec.Codec}
    \lean{ProtocolSpec.Codec.mk'}
    \lean{ProtocolSpec.Codec.instSerializeMessage}
    \lean{ProtocolSpec.Codec.instDeserializeChallenge}
\end{definition}

\begin{definition}[Salt Codec]
    \label{def:duplex_sponge_salt_codec}
    The paper-facing salted transform distinguishes the on-sponge salt
    $\tau \in \Sigma^\delta$ from the pre-encoded salt value used by the canonical FS-standard
    hybrids. ArkLib packages this bridge as a \emph{salt codec}, with an injective encoding and a
    left inverse.
    \lean{SaltCodec}
    \lean{SaltCodec.encode_injective}
\end{definition}

\subsection{Duplex-Sponge Fiat-Shamir Transform}\label{sec:fiat_shamir_duplex_transform}

\begin{definition}[Duplex-Sponge Transcript Derivation]
    \label{def:derive_transcript_dsfs}
    Given a statement and a non-interactive proof containing all prover messages, the verifier
    starts a canonical duplex sponge at the statement, absorbs each encoded prover message, and
    squeezes and decodes each verifier challenge. The result is the same kind of full transcript
    expected by the original interactive verifier.
    \lean{ProtocolSpec.Messages.deriveTranscriptDSFS}
    \lean{ProtocolSpec.Messages.deriveTranscriptDSFSSalted}
    \uses{def:duplex_sponge_challenge_oracle, def:duplex_sponge_codec}
\end{definition}

\begin{definition}[Duplex-Sponge Fiat-Shamir Reduction]
    \label{def:duplex_sponge_fiat_shamir}
    The duplex-sponge Fiat-Shamir transformation applies the sponge-aware prover and verifier
    transformations to an interactive reduction:
    \[
        R.\mathsf{duplexSpongeFiatShamir}
        : \NonInteractiveReduction
          \;(\forall i,\pSpec.\Message(i))\;
          (\oSpec + \mathcal{D}_{\mathfrak{S}})\;\cdots .
    \]
    The prover performs the same absorb/squeeze process while computing its messages; the verifier
    independently replays the sponge transcript derivation.
    \lean{Prover.processRoundDSFS}
    \lean{Prover.runToRoundDSFS}
    \lean{Prover.duplexSpongeFiatShamir}
    \lean{Verifier.duplexSpongeFiatShamir}
    \lean{Reduction.duplexSpongeFiatShamir}
    \uses{def:derive_transcript_dsfs}
\end{definition}

\begin{definition}[Salted Duplex-Sponge Fiat-Shamir]
    \label{def:salted_duplex_sponge_fiat_shamir}
    The salted wrapper follows the paper-native construction by sampling or supplying a salt
    $\tau \in \Sigma^\delta$, absorbing it once after sponge initialization, and then running the
    same message/challenge absorb-squeeze loop. ArkLib exposes both an explicit salt-source
    wrapper and a uniform-salt convenience wrapper when the ambient oracle context contains the
    uniform sampler.
    \lean{uniformSalt}
    \lean{Prover.runToRoundDSFSSalted}
    \lean{Prover.duplexSpongeFiatShamirSalted}
    \lean{Prover.duplexSpongeFiatShamirSaltedRandom}
    \lean{Verifier.duplexSpongeFiatShamirSalted}
    \lean{Reduction.duplexSpongeFiatShamirSalted}
    \lean{Reduction.duplexSpongeFiatShamirSaltedRandom}
    \uses{def:duplex_sponge_salt_codec, def:duplex_sponge_fiat_shamir}
\end{definition}

\subsection{Security Status}\label{sec:fiat_shamir_security}

\begin{theorem}[Duplex-Sponge Completeness Unroll]
    \label{thm:duplex_sponge_completeness_unroll}
    The DSFS run-collapse obligations for both the unsalted and salted transforms are discharged.
    Consequently, completeness of the transformed reduction is unconditionally equivalent to the
    explicit honest DSFS execution packaged in the security layer.
    \lean{Reduction.duplexSpongeFiatShamir_runCollapse}
    \lean{Reduction.duplexSpongeFiatShamirSalted_runCollapse}
    \lean{Reduction.duplexSpongeFiatShamir_completeness_unroll_discharged}
    \lean{Reduction.duplexSpongeFiatShamirSalted_completeness_unroll_discharged}
    \uses{def:duplex_sponge_fiat_shamir, def:salted_duplex_sponge_fiat_shamir}
\end{theorem}

\begin{definition}[DSFS Key Lemma Statement]
    \label{def:duplex_sponge_key_lemma}
    The DSFS security framework compares the duplex-sponge game against the canonical
    Fiat-Shamir game after remapping the DSFS query traces to the canonical challenge-oracle
    surface. The core statement is a total-variation bound: for every query-bounded malicious
    DSFS prover, there exists a simulated canonical-FS prover whose game output distribution is
    within the error term $\eta^\star$ of the remapped DSFS game.
    \lean{DuplexSpongeFS.duplexSpongeFiatShamirGame}
    \lean{DuplexSpongeFS.duplexSpongeFiatShamirGameRemapped}
    \lean{DuplexSpongeFS.KeyLemmaStatement}
    \lean{DuplexSpongeFS.KeyLemmaResidual}
    \lean{DuplexSpongeFS.duplexSpongeToFSGameStatDist}
    \uses{def:fiat_shamir_challenge_oracle_interface, def:duplex_sponge_challenge_oracle}
\end{definition}

\begin{remark}[What Is Proved Versus Residualized]
    The executable duplex-sponge specification, codec layer, salted variants, trace/prover
    transformation infrastructure, run-collapse completeness bridge, and the canonical
    Fiat-Shamir transfer theorems are machine checked.

    The full CO25 Lemma 5.1 hybrid proof for DSFS remains an explicit post-grant security
    hardening assumption, not a hidden proof hole. The strict residual surface is:
    \[
    \begin{gathered}
    \mathsf{DuplexSpongeFS.KeyLemmaResidual},\\
    \mathsf{Lemma5\_12HonestResidual},\quad
    \mathsf{Lemma5\_14HonestResidual},\quad
    \mathsf{Lemma5\_16HonestResidual},\\
    \mathsf{SimulatedProverChallengeBudgetResidual},\quad
    \mathsf{SimulatedProverSharedBudgetResidual},\\
    \mathsf{KeyLemmaEagerResidual},\\
    \mathsf{D2sQueryStepGSpecBudgetResidual},\quad
    \mathsf{D2fOuterImplSharedBudgetResidual}.
    \end{gathered}
    \]
    These residuals express the remaining CO25 Section 5 hybrid-game proof obligations:
    dispatcher query budgets, honest bad-event implications, simulated-prover budgets, and the
    final eager-oracle key lemma. They are deliberately exposed as named propositions consumed by
    downstream DSFS soundness reductions.
    \lean{DuplexSpongeFS.KeyLemmaFoundations.KeyLemmaStatementEager}
    \lean{DuplexSpongeFS.KeyLemmaFoundations.KeyLemmaEagerResidual}
    \lean{DuplexSpongeFS.KeyLemmaFoundations.D2sQueryStepGSpecBudgetResidual}
    \lean{DuplexSpongeFS.KeyLemmaFoundations.D2fOuterImplSharedBudgetResidual}
\end{remark}
